\section{System Design}
\label{sec:design}

\subsection{Architecture}
\label{sec:architecture}

The system is a Python package, \term{vidbpedia}, organised in four layers
(Figure~\ref{fig:architecture}). Each step is one command,
\term{python -m vidbpedia <step>}: \term{seeds}, \term{enrich},
\term{ontology}, \term{build}, \term{postprocess} and \term{serve}, plus
\term{query} for SPARQL from the terminal. Only the collection layer uses the
network: HTTP responses are cached in SQLite and the raw downloads are
committed under \term{data/raw/}, so the later layers rebuild the dataset
offline and identically. Graphs are built with the rdflib API~\cite{rdflib}
rather than by concatenating Turtle strings, which rules out syntax and
escaping errors, and a single function creates every resource IRI, so an
entity cannot receive two IRIs.

\begin{figure}[t]
\centering
\begin{adjustbox}{max width=\textwidth}
\begin{tikzpicture}[
    font=\scriptsize, >={Stealth[length=1.8mm]},
    box/.style={draw=accgray, rounded corners=2pt, align=center, text width=2.7cm,
                minimum height=1.3cm, inner sep=3pt},
    stage/.style={box, draw=accblue, fill=accblue!7},
    store/.style={box, fill=accgray!12},
    ext/.style={box, dashed},
    flow/.style={->, accgray, thick},
  ]
  \node[ext]                        (src)   {Wikidata Query Service\\[1pt] MediaWiki API\\ (vi.wikipedia.org)};
  \node[stage, right=0.45cm of src]   (crawl) {\textbf{(1) Collection}\\[1pt] \texttt{seeds}, \texttt{enrich}\\ SQLite HTTP cache};
  \node[store, right=0.45cm of crawl] (raw)   {\texttt{data/raw/} (committed)\\[1pt] Wikidata facts,\\ article metadata, infoboxes};
  \node[stage, right=0.45cm of raw]   (build) {\textbf{(2) Semantics}\\[1pt] \texttt{ontology}: merge modules\\ \texttt{build}: infobox $\rightarrow$ RDF};
  \node[stage, below=0.75cm of build] (post)  {\textbf{(3) Post-processing}\\[1pt] validate $\rightarrow$ OWL\,2\,RL\\ $\rightarrow$ VoID $\rightarrow$ export};
  \node[store, left=0.45cm of post]   (out)   {\texttt{data/}\\[1pt] 131{,}543 triples (.ttl, .nt)\\ ontology, asserted\\ and inferred parts};
  \node[stage, left=0.45cm of out]    (srv)   {\textbf{(4) Server}\\[1pt] FastAPI + React UI\\ in-memory rdflib graph\\ \texttt{/sparql} endpoint};
  \node[ext, left=0.45cm of srv]      (users) {Browsers, terminal,\\ SPARQL and RDF clients\\[1pt] LLM (OpenAI-compatible\\ API, optional)};
  \draw[flow] (src) -- (crawl);
  \draw[flow] (crawl) -- (raw);
  \draw[flow] (raw) -- (build);
  \draw[flow] (build) -- node[right, font=\tiny, text=accgray] {\texttt{data/raw/rdf/*.ttl}} (post);
  \draw[flow] (post) -- (out);
  \draw[flow] (out) -- (srv);
  \draw[flow, <->] (srv) -- (users);
\end{tikzpicture}
\end{adjustbox}
\caption{Architecture. Blue boxes are the four layers of the \texttt{vidbpedia}
package (the interface of layer~(4) is in \texttt{ui/}), gray boxes are files on
disk, dashed boxes are external systems. Only layer~(1) calls the network.}
\label{fig:architecture}
\end{figure}

% =========================================================================
\subsection{Ontology}
\label{sec:ontology}

The ontology is written as ten Turtle modules in \term{ontology/}, from
\term{000-prefixes.ttl} to \term{420-football.ttl}. They are merged in numeric order into
\term{vi-ontology.ttl} (859 triples) and checked by rdflib when merged.
Table~\ref{tab:namespaces} lists the namespaces. Following
DBpedia~\cite{lehmann2015dbpedia}, resources are named by their article title,
infobox parameters have a raw namespace, and categories are SKOS concepts.

\begin{table}[t]
\centering
\caption{Namespaces. The \texttt{vi.dbpedia.org} namespaces mirror the IRI
scheme of DBpedia language chapters.}
\label{tab:namespaces}
\small
\begin{tabularx}{\textwidth}{@{}ll>{\raggedright\arraybackslash}X@{}}
\toprule
Prefix & Namespace & Used for \\
\midrule
\term{vio:}  & \term{http://vi.dbpedia.org/ontology/} & classes and properties of this project \\
\term{vres:} & \term{http://vi.dbpedia.org/resource/} & resources (entities) \\
\term{vip:}  & \term{http://vi.dbpedia.org/property/} & raw infobox properties, like \term{dbp:} \\
\term{vcat:} & \term{http://vi.dbpedia.org/resource/Thể\_loại:} & Wikipedia categories \\
\term{dbo:}  & \term{http://dbpedia.org/ontology/} & DBpedia Ontology \\
\term{dbr:}  & \term{http://dbpedia.org/resource/} & English DBpedia resources \\
\term{wd:}   & \term{http://www.wikidata.org/entity/} & Wikidata items \\
\bottomrule
\end{tabularx}
\end{table}

\paragraph{Classes.} The ontology has two layers~\cite{heath2011linked}. The
upper layer reuses the DBpedia Ontology: \term{050-dbo-alignment.ttl} declares
the 24 DBpedia classes and 44 properties that \term{vio:} maps to, with
DBpedia's English labels, our Vietnamese labels (\term{dbo:Animal} is ``Động
vật'') and \term{rdfs:isDefinedBy dbo:}, but not DBpedia's domains and ranges.
The lower layer has 18 \term{vio:} classes in four disjoint trees: people,
organisations, locations and career stations (Table~\ref{tab:classes}). A
\term{vio:} class exists only where an axiom or a competency
question~\cite{gruninger1995methodology} needs it, such as
\term{vio:FormerProvince}; there is no \term{vio:Animal}. Each class is an
\term{rdfs:subClassOf}, not an \term{owl:equivalentClass}, of a DBpedia class,
directly or through its parent. Data described with \term{vio:} can thus be
queried with \term{dbo:} after reasoning (Section~\ref{sec:reasoning}), while
our domains, functional properties and disjointness axioms constrain only our
own data, avoiding the Exclusivity anti-pattern~\cite{paulheim2024modeling}.
\term{vio:FormerProvince} holds the
provinces that were dissolved or merged, including in the 2025 reorganisation:
a province is former if Wikidata gives a dissolution date (\term{P576}) or no
current province statement (\term{P31}). With this rule, ``how many provinces are there now'' returns
34, while history is kept through \term{vio:successor} and
\term{vio:predecessor}.

\begin{figure}[t]
\centering
\resizebox{\textwidth}{!}{% RDFS graph of the vio: ontology in the lecture style: ellipses for resources,
% rdfs:Class / rdf:Property at the top, the schema above a horizontal line and
% instance data below it, rdf:type edges crossing the line.
% An excerpt: the full class hierarchy is in the report's class table.
%
% Shared by report/ (% RDFS graph of the vio: ontology in the lecture style: ellipses for resources,
% rdfs:Class / rdf:Property at the top, the schema above a horizontal line and
% instance data below it, rdf:type edges crossing the line.
% An excerpt: the full class hierarchy is in the report's class table.
%
% Shared by report/ (% RDFS graph of the vio: ontology in the lecture style: ellipses for resources,
% rdfs:Class / rdf:Property at the top, the schema above a horizontal line and
% instance data below it, rdf:type edges crossing the line.
% An excerpt: the full class hierarchy is in the report's class table.
%
% Shared by report/ (\input{figures/rdfs_graph}) and slides/
% (\input{../report/figures/rdfs_graph}). The including preamble defines the
% colours rdfsdraw, rdfsfill (instances), rdfsprop (properties) and rdfsaxiom
% (entailed triples).
\begin{tikzpicture}[
    % absolute font sizes, so the figure looks the same in the report (10pt) and in beamer (11pt)
    font=\fontsize{7}{8.5}\selectfont, >={Stealth[length=1.7mm]},
    res/.style={ellipse, draw=black!80, fill=white, inner sep=1pt, minimum height=0.72cm,
                font=\fontsize{7}{8.5}\selectfont\ttfamily},
    meta/.style={res, fill=black!8, draw=black!60},
    dbo/.style={res, dashed, draw=black!55, text=black!70},
    prp/.style={res, draw=rdfsprop, text=rdfsprop},
    dboprp/.style={prp, dashed},
    inst/.style={res, fill=rdfsfill, draw=rdfsdraw},
    e/.style={->, black!80},
    typ/.style={->, black!45},
    dr/.style={->, rdfsprop},
    lbl/.style={font=\fontsize{5}{6}\selectfont\itshape, fill=white, inner sep=0.8pt},
  ]
  % ---- RDFS vocabulary ----
  \node[meta] (class)    at (4.7,1.75)  {rdfs:Class};
  \node[meta] (property) at (10.6,1.75) {rdf:Property};

  % ---- classes ----
  \node[dbo] (dboanimal) at (1.8,1.75)  {dbo:Animal};
  \node[dbo] (dboperson) at (1.8,0.35)  {dbo:Person};
  \node[res] (person)    at (1.8,-1.05) {vio:Person};
  \node[res] (athlete)   at (1.8,-2.45) {vio:Athlete};
  \node[res] (fp)        at (1.8,-3.85) {vio:FootballPlayer};
  \node[res] (cs)        at (6.2,-2.45) {vio:CareerStation};
  \node[res] (club)      at (6.2,-3.85) {vio:ClubStation};
  \node[res] (org)       at (10.4,-2.45) {vio:Organisation};
  \node[res] (fc)        at (10.4,-3.85) {vio:FootballClub};
  \node[res] (loc)       at (14.7,-2.45) {vio:Location};
  \node[res] (prov)      at (14.7,-3.85) {vio:Province};

  % ---- properties ----
  \node[prp]    (cstation)  at (6.2,0.35)   {vio:careerStation};
  \node[prp]    (bplace)    at (10.6,0.35)  {vio:birthPlace};
  \node[dboprp] (dbobplace) at (15.0,0.35)  {dbo:birthPlace};
  \node[prp]    (bprov)     at (10.6,-1.05) {vio:birthProvince};

  % ---- instance data ----
  \node[inst] (cp)   at (1.8,-6.0)  {vres:Nguyễn\_Công\_Phượng};
  \node[inst] (st2)  at (6.2,-6.0)  {vres:…\_\_2};
  \node[inst] (hagl) at (10.4,-6.0) {vres:…Hoàng\_Anh\_Gia\_Lai};
  \node[inst] (na)   at (14.7,-6.0) {vres:Nghệ\_An};

  % ---- schema / data line ----
  \draw[black!45, thick] (-0.9,-4.95) -- (16.7,-4.95);
  \node[lbl, fill=none, text=black!55, anchor=south west, inner sep=2pt] at (-0.9,-4.95) {schema (RDFS)};
  \node[lbl, fill=none, text=black!55, anchor=north west, inner sep=2pt] at (-0.9,-4.95) {data (RDF)};

  % ---- rdfs:subClassOf ----
  \draw[e] (dboperson) -- node[lbl, right] {rdfs:subClassOf} (dboanimal);
  \draw[e] (person)  -- node[lbl, right] {rdfs:subClassOf} (dboperson);
  \draw[e] (athlete) -- node[lbl, right] {rdfs:subClassOf} (person);
  \draw[e] (fp)      -- node[lbl, right] {rdfs:subClassOf} (athlete);
  \draw[e] (club)    -- node[lbl, right] {rdfs:subClassOf} (cs);
  \draw[e] (fc)      -- node[lbl, right] {rdfs:subClassOf} (org);
  \draw[e] (prov)    -- node[lbl, right] {rdfs:subClassOf} (loc);

  % ---- rdfs:subPropertyOf ----
  \draw[e] (bprov)  -- node[lbl, right] {rdfs:subPropertyOf} (bplace);
  \draw[e] (bplace) -- node[lbl, above] {rdfs:subPropertyOf} (dbobplace);

  % ---- rdf:type to the RDFS vocabulary ----
  \draw[typ] (person.north east) -- node[lbl, sloped, above] {rdf:type} (class.south west);
  \draw[typ] (cstation) -- node[lbl, sloped, above] {rdf:type} (property);
  \draw[typ] (bplace) -- node[lbl, right] {rdf:type} (property);
  \draw[typ] (dbobplace) -- node[lbl, sloped, above] {rdf:type} (property);

  % ---- rdfs:domain / rdfs:range ----
  \draw[dr] (cstation.south west) -- node[lbl, sloped, above, text=rdfsprop] {rdfs:domain} (fp.north east);
  \draw[dr] (cstation) -- node[lbl, right, text=rdfsprop, pos=0.45] {rdfs:range} (cs);
  \draw[dr] (bprov.south east) -- node[lbl, sloped, above, text=rdfsprop] {rdfs:range} (prov.north west);

  % ---- instance triples ----
  \draw[e] (cp) -- node[lbl, above] {vio:careerStation} (st2);
  \draw[e] (st2) -- node[lbl, above] {vio:team} (hagl);
  \draw[e] (cp.south) -- ++(0,-0.55) -| node[lbl, pos=0.25] {vio:birthProvince} (na.south);

  % ---- rdf:type across the line ----
  \draw[typ] (cp)   -- node[lbl, right] {rdf:type} (fp);
  \draw[typ] (st2)  -- node[lbl, right] {rdf:type} (club);
  \draw[typ] (hagl) -- node[lbl, right] {rdf:type} (fc);
  \draw[typ] (na)   -- node[lbl, right] {rdf:type} (prov);

  % ---- a triple entailed by RDFS (rdfs9: type propagates up subClassOf) ----
  \draw[->, rdfsaxiom, dashed, thick] (cp.west) to[out=170, in=190, looseness=1.15]
      node[lbl, text=rdfsaxiom, sloped, above, pos=0.5] {rdf:type (entailed)} (dboperson.west);
\end{tikzpicture}
) and slides/
% (% RDFS graph of the vio: ontology in the lecture style: ellipses for resources,
% rdfs:Class / rdf:Property at the top, the schema above a horizontal line and
% instance data below it, rdf:type edges crossing the line.
% An excerpt: the full class hierarchy is in the report's class table.
%
% Shared by report/ (\input{figures/rdfs_graph}) and slides/
% (\input{../report/figures/rdfs_graph}). The including preamble defines the
% colours rdfsdraw, rdfsfill (instances), rdfsprop (properties) and rdfsaxiom
% (entailed triples).
\begin{tikzpicture}[
    % absolute font sizes, so the figure looks the same in the report (10pt) and in beamer (11pt)
    font=\fontsize{7}{8.5}\selectfont, >={Stealth[length=1.7mm]},
    res/.style={ellipse, draw=black!80, fill=white, inner sep=1pt, minimum height=0.72cm,
                font=\fontsize{7}{8.5}\selectfont\ttfamily},
    meta/.style={res, fill=black!8, draw=black!60},
    dbo/.style={res, dashed, draw=black!55, text=black!70},
    prp/.style={res, draw=rdfsprop, text=rdfsprop},
    dboprp/.style={prp, dashed},
    inst/.style={res, fill=rdfsfill, draw=rdfsdraw},
    e/.style={->, black!80},
    typ/.style={->, black!45},
    dr/.style={->, rdfsprop},
    lbl/.style={font=\fontsize{5}{6}\selectfont\itshape, fill=white, inner sep=0.8pt},
  ]
  % ---- RDFS vocabulary ----
  \node[meta] (class)    at (4.7,1.75)  {rdfs:Class};
  \node[meta] (property) at (10.6,1.75) {rdf:Property};

  % ---- classes ----
  \node[dbo] (dboanimal) at (1.8,1.75)  {dbo:Animal};
  \node[dbo] (dboperson) at (1.8,0.35)  {dbo:Person};
  \node[res] (person)    at (1.8,-1.05) {vio:Person};
  \node[res] (athlete)   at (1.8,-2.45) {vio:Athlete};
  \node[res] (fp)        at (1.8,-3.85) {vio:FootballPlayer};
  \node[res] (cs)        at (6.2,-2.45) {vio:CareerStation};
  \node[res] (club)      at (6.2,-3.85) {vio:ClubStation};
  \node[res] (org)       at (10.4,-2.45) {vio:Organisation};
  \node[res] (fc)        at (10.4,-3.85) {vio:FootballClub};
  \node[res] (loc)       at (14.7,-2.45) {vio:Location};
  \node[res] (prov)      at (14.7,-3.85) {vio:Province};

  % ---- properties ----
  \node[prp]    (cstation)  at (6.2,0.35)   {vio:careerStation};
  \node[prp]    (bplace)    at (10.6,0.35)  {vio:birthPlace};
  \node[dboprp] (dbobplace) at (15.0,0.35)  {dbo:birthPlace};
  \node[prp]    (bprov)     at (10.6,-1.05) {vio:birthProvince};

  % ---- instance data ----
  \node[inst] (cp)   at (1.8,-6.0)  {vres:Nguyễn\_Công\_Phượng};
  \node[inst] (st2)  at (6.2,-6.0)  {vres:…\_\_2};
  \node[inst] (hagl) at (10.4,-6.0) {vres:…Hoàng\_Anh\_Gia\_Lai};
  \node[inst] (na)   at (14.7,-6.0) {vres:Nghệ\_An};

  % ---- schema / data line ----
  \draw[black!45, thick] (-0.9,-4.95) -- (16.7,-4.95);
  \node[lbl, fill=none, text=black!55, anchor=south west, inner sep=2pt] at (-0.9,-4.95) {schema (RDFS)};
  \node[lbl, fill=none, text=black!55, anchor=north west, inner sep=2pt] at (-0.9,-4.95) {data (RDF)};

  % ---- rdfs:subClassOf ----
  \draw[e] (dboperson) -- node[lbl, right] {rdfs:subClassOf} (dboanimal);
  \draw[e] (person)  -- node[lbl, right] {rdfs:subClassOf} (dboperson);
  \draw[e] (athlete) -- node[lbl, right] {rdfs:subClassOf} (person);
  \draw[e] (fp)      -- node[lbl, right] {rdfs:subClassOf} (athlete);
  \draw[e] (club)    -- node[lbl, right] {rdfs:subClassOf} (cs);
  \draw[e] (fc)      -- node[lbl, right] {rdfs:subClassOf} (org);
  \draw[e] (prov)    -- node[lbl, right] {rdfs:subClassOf} (loc);

  % ---- rdfs:subPropertyOf ----
  \draw[e] (bprov)  -- node[lbl, right] {rdfs:subPropertyOf} (bplace);
  \draw[e] (bplace) -- node[lbl, above] {rdfs:subPropertyOf} (dbobplace);

  % ---- rdf:type to the RDFS vocabulary ----
  \draw[typ] (person.north east) -- node[lbl, sloped, above] {rdf:type} (class.south west);
  \draw[typ] (cstation) -- node[lbl, sloped, above] {rdf:type} (property);
  \draw[typ] (bplace) -- node[lbl, right] {rdf:type} (property);
  \draw[typ] (dbobplace) -- node[lbl, sloped, above] {rdf:type} (property);

  % ---- rdfs:domain / rdfs:range ----
  \draw[dr] (cstation.south west) -- node[lbl, sloped, above, text=rdfsprop] {rdfs:domain} (fp.north east);
  \draw[dr] (cstation) -- node[lbl, right, text=rdfsprop, pos=0.45] {rdfs:range} (cs);
  \draw[dr] (bprov.south east) -- node[lbl, sloped, above, text=rdfsprop] {rdfs:range} (prov.north west);

  % ---- instance triples ----
  \draw[e] (cp) -- node[lbl, above] {vio:careerStation} (st2);
  \draw[e] (st2) -- node[lbl, above] {vio:team} (hagl);
  \draw[e] (cp.south) -- ++(0,-0.55) -| node[lbl, pos=0.25] {vio:birthProvince} (na.south);

  % ---- rdf:type across the line ----
  \draw[typ] (cp)   -- node[lbl, right] {rdf:type} (fp);
  \draw[typ] (st2)  -- node[lbl, right] {rdf:type} (club);
  \draw[typ] (hagl) -- node[lbl, right] {rdf:type} (fc);
  \draw[typ] (na)   -- node[lbl, right] {rdf:type} (prov);

  % ---- a triple entailed by RDFS (rdfs9: type propagates up subClassOf) ----
  \draw[->, rdfsaxiom, dashed, thick] (cp.west) to[out=170, in=190, looseness=1.15]
      node[lbl, text=rdfsaxiom, sloped, above, pos=0.5] {rdf:type (entailed)} (dboperson.west);
\end{tikzpicture}
). The including preamble defines the
% colours rdfsdraw, rdfsfill (instances), rdfsprop (properties) and rdfsaxiom
% (entailed triples).
\begin{tikzpicture}[
    % absolute font sizes, so the figure looks the same in the report (10pt) and in beamer (11pt)
    font=\fontsize{7}{8.5}\selectfont, >={Stealth[length=1.7mm]},
    res/.style={ellipse, draw=black!80, fill=white, inner sep=1pt, minimum height=0.72cm,
                font=\fontsize{7}{8.5}\selectfont\ttfamily},
    meta/.style={res, fill=black!8, draw=black!60},
    dbo/.style={res, dashed, draw=black!55, text=black!70},
    prp/.style={res, draw=rdfsprop, text=rdfsprop},
    dboprp/.style={prp, dashed},
    inst/.style={res, fill=rdfsfill, draw=rdfsdraw},
    e/.style={->, black!80},
    typ/.style={->, black!45},
    dr/.style={->, rdfsprop},
    lbl/.style={font=\fontsize{5}{6}\selectfont\itshape, fill=white, inner sep=0.8pt},
  ]
  % ---- RDFS vocabulary ----
  \node[meta] (class)    at (4.7,1.75)  {rdfs:Class};
  \node[meta] (property) at (10.6,1.75) {rdf:Property};

  % ---- classes ----
  \node[dbo] (dboanimal) at (1.8,1.75)  {dbo:Animal};
  \node[dbo] (dboperson) at (1.8,0.35)  {dbo:Person};
  \node[res] (person)    at (1.8,-1.05) {vio:Person};
  \node[res] (athlete)   at (1.8,-2.45) {vio:Athlete};
  \node[res] (fp)        at (1.8,-3.85) {vio:FootballPlayer};
  \node[res] (cs)        at (6.2,-2.45) {vio:CareerStation};
  \node[res] (club)      at (6.2,-3.85) {vio:ClubStation};
  \node[res] (org)       at (10.4,-2.45) {vio:Organisation};
  \node[res] (fc)        at (10.4,-3.85) {vio:FootballClub};
  \node[res] (loc)       at (14.7,-2.45) {vio:Location};
  \node[res] (prov)      at (14.7,-3.85) {vio:Province};

  % ---- properties ----
  \node[prp]    (cstation)  at (6.2,0.35)   {vio:careerStation};
  \node[prp]    (bplace)    at (10.6,0.35)  {vio:birthPlace};
  \node[dboprp] (dbobplace) at (15.0,0.35)  {dbo:birthPlace};
  \node[prp]    (bprov)     at (10.6,-1.05) {vio:birthProvince};

  % ---- instance data ----
  \node[inst] (cp)   at (1.8,-6.0)  {vres:Nguyễn\_Công\_Phượng};
  \node[inst] (st2)  at (6.2,-6.0)  {vres:…\_\_2};
  \node[inst] (hagl) at (10.4,-6.0) {vres:…Hoàng\_Anh\_Gia\_Lai};
  \node[inst] (na)   at (14.7,-6.0) {vres:Nghệ\_An};

  % ---- schema / data line ----
  \draw[black!45, thick] (-0.9,-4.95) -- (16.7,-4.95);
  \node[lbl, fill=none, text=black!55, anchor=south west, inner sep=2pt] at (-0.9,-4.95) {schema (RDFS)};
  \node[lbl, fill=none, text=black!55, anchor=north west, inner sep=2pt] at (-0.9,-4.95) {data (RDF)};

  % ---- rdfs:subClassOf ----
  \draw[e] (dboperson) -- node[lbl, right] {rdfs:subClassOf} (dboanimal);
  \draw[e] (person)  -- node[lbl, right] {rdfs:subClassOf} (dboperson);
  \draw[e] (athlete) -- node[lbl, right] {rdfs:subClassOf} (person);
  \draw[e] (fp)      -- node[lbl, right] {rdfs:subClassOf} (athlete);
  \draw[e] (club)    -- node[lbl, right] {rdfs:subClassOf} (cs);
  \draw[e] (fc)      -- node[lbl, right] {rdfs:subClassOf} (org);
  \draw[e] (prov)    -- node[lbl, right] {rdfs:subClassOf} (loc);

  % ---- rdfs:subPropertyOf ----
  \draw[e] (bprov)  -- node[lbl, right] {rdfs:subPropertyOf} (bplace);
  \draw[e] (bplace) -- node[lbl, above] {rdfs:subPropertyOf} (dbobplace);

  % ---- rdf:type to the RDFS vocabulary ----
  \draw[typ] (person.north east) -- node[lbl, sloped, above] {rdf:type} (class.south west);
  \draw[typ] (cstation) -- node[lbl, sloped, above] {rdf:type} (property);
  \draw[typ] (bplace) -- node[lbl, right] {rdf:type} (property);
  \draw[typ] (dbobplace) -- node[lbl, sloped, above] {rdf:type} (property);

  % ---- rdfs:domain / rdfs:range ----
  \draw[dr] (cstation.south west) -- node[lbl, sloped, above, text=rdfsprop] {rdfs:domain} (fp.north east);
  \draw[dr] (cstation) -- node[lbl, right, text=rdfsprop, pos=0.45] {rdfs:range} (cs);
  \draw[dr] (bprov.south east) -- node[lbl, sloped, above, text=rdfsprop] {rdfs:range} (prov.north west);

  % ---- instance triples ----
  \draw[e] (cp) -- node[lbl, above] {vio:careerStation} (st2);
  \draw[e] (st2) -- node[lbl, above] {vio:team} (hagl);
  \draw[e] (cp.south) -- ++(0,-0.55) -| node[lbl, pos=0.25] {vio:birthProvince} (na.south);

  % ---- rdf:type across the line ----
  \draw[typ] (cp)   -- node[lbl, right] {rdf:type} (fp);
  \draw[typ] (st2)  -- node[lbl, right] {rdf:type} (club);
  \draw[typ] (hagl) -- node[lbl, right] {rdf:type} (fc);
  \draw[typ] (na)   -- node[lbl, right] {rdf:type} (prov);

  % ---- a triple entailed by RDFS (rdfs9: type propagates up subClassOf) ----
  \draw[->, rdfsaxiom, dashed, thick] (cp.west) to[out=170, in=190, looseness=1.15]
      node[lbl, text=rdfsaxiom, sloped, above, pos=0.5] {rdf:type (entailed)} (dboperson.west);
\end{tikzpicture}
) and slides/
% (% RDFS graph of the vio: ontology in the lecture style: ellipses for resources,
% rdfs:Class / rdf:Property at the top, the schema above a horizontal line and
% instance data below it, rdf:type edges crossing the line.
% An excerpt: the full class hierarchy is in the report's class table.
%
% Shared by report/ (% RDFS graph of the vio: ontology in the lecture style: ellipses for resources,
% rdfs:Class / rdf:Property at the top, the schema above a horizontal line and
% instance data below it, rdf:type edges crossing the line.
% An excerpt: the full class hierarchy is in the report's class table.
%
% Shared by report/ (\input{figures/rdfs_graph}) and slides/
% (\input{../report/figures/rdfs_graph}). The including preamble defines the
% colours rdfsdraw, rdfsfill (instances), rdfsprop (properties) and rdfsaxiom
% (entailed triples).
\begin{tikzpicture}[
    % absolute font sizes, so the figure looks the same in the report (10pt) and in beamer (11pt)
    font=\fontsize{7}{8.5}\selectfont, >={Stealth[length=1.7mm]},
    res/.style={ellipse, draw=black!80, fill=white, inner sep=1pt, minimum height=0.72cm,
                font=\fontsize{7}{8.5}\selectfont\ttfamily},
    meta/.style={res, fill=black!8, draw=black!60},
    dbo/.style={res, dashed, draw=black!55, text=black!70},
    prp/.style={res, draw=rdfsprop, text=rdfsprop},
    dboprp/.style={prp, dashed},
    inst/.style={res, fill=rdfsfill, draw=rdfsdraw},
    e/.style={->, black!80},
    typ/.style={->, black!45},
    dr/.style={->, rdfsprop},
    lbl/.style={font=\fontsize{5}{6}\selectfont\itshape, fill=white, inner sep=0.8pt},
  ]
  % ---- RDFS vocabulary ----
  \node[meta] (class)    at (4.7,1.75)  {rdfs:Class};
  \node[meta] (property) at (10.6,1.75) {rdf:Property};

  % ---- classes ----
  \node[dbo] (dboanimal) at (1.8,1.75)  {dbo:Animal};
  \node[dbo] (dboperson) at (1.8,0.35)  {dbo:Person};
  \node[res] (person)    at (1.8,-1.05) {vio:Person};
  \node[res] (athlete)   at (1.8,-2.45) {vio:Athlete};
  \node[res] (fp)        at (1.8,-3.85) {vio:FootballPlayer};
  \node[res] (cs)        at (6.2,-2.45) {vio:CareerStation};
  \node[res] (club)      at (6.2,-3.85) {vio:ClubStation};
  \node[res] (org)       at (10.4,-2.45) {vio:Organisation};
  \node[res] (fc)        at (10.4,-3.85) {vio:FootballClub};
  \node[res] (loc)       at (14.7,-2.45) {vio:Location};
  \node[res] (prov)      at (14.7,-3.85) {vio:Province};

  % ---- properties ----
  \node[prp]    (cstation)  at (6.2,0.35)   {vio:careerStation};
  \node[prp]    (bplace)    at (10.6,0.35)  {vio:birthPlace};
  \node[dboprp] (dbobplace) at (15.0,0.35)  {dbo:birthPlace};
  \node[prp]    (bprov)     at (10.6,-1.05) {vio:birthProvince};

  % ---- instance data ----
  \node[inst] (cp)   at (1.8,-6.0)  {vres:Nguyễn\_Công\_Phượng};
  \node[inst] (st2)  at (6.2,-6.0)  {vres:…\_\_2};
  \node[inst] (hagl) at (10.4,-6.0) {vres:…Hoàng\_Anh\_Gia\_Lai};
  \node[inst] (na)   at (14.7,-6.0) {vres:Nghệ\_An};

  % ---- schema / data line ----
  \draw[black!45, thick] (-0.9,-4.95) -- (16.7,-4.95);
  \node[lbl, fill=none, text=black!55, anchor=south west, inner sep=2pt] at (-0.9,-4.95) {schema (RDFS)};
  \node[lbl, fill=none, text=black!55, anchor=north west, inner sep=2pt] at (-0.9,-4.95) {data (RDF)};

  % ---- rdfs:subClassOf ----
  \draw[e] (dboperson) -- node[lbl, right] {rdfs:subClassOf} (dboanimal);
  \draw[e] (person)  -- node[lbl, right] {rdfs:subClassOf} (dboperson);
  \draw[e] (athlete) -- node[lbl, right] {rdfs:subClassOf} (person);
  \draw[e] (fp)      -- node[lbl, right] {rdfs:subClassOf} (athlete);
  \draw[e] (club)    -- node[lbl, right] {rdfs:subClassOf} (cs);
  \draw[e] (fc)      -- node[lbl, right] {rdfs:subClassOf} (org);
  \draw[e] (prov)    -- node[lbl, right] {rdfs:subClassOf} (loc);

  % ---- rdfs:subPropertyOf ----
  \draw[e] (bprov)  -- node[lbl, right] {rdfs:subPropertyOf} (bplace);
  \draw[e] (bplace) -- node[lbl, above] {rdfs:subPropertyOf} (dbobplace);

  % ---- rdf:type to the RDFS vocabulary ----
  \draw[typ] (person.north east) -- node[lbl, sloped, above] {rdf:type} (class.south west);
  \draw[typ] (cstation) -- node[lbl, sloped, above] {rdf:type} (property);
  \draw[typ] (bplace) -- node[lbl, right] {rdf:type} (property);
  \draw[typ] (dbobplace) -- node[lbl, sloped, above] {rdf:type} (property);

  % ---- rdfs:domain / rdfs:range ----
  \draw[dr] (cstation.south west) -- node[lbl, sloped, above, text=rdfsprop] {rdfs:domain} (fp.north east);
  \draw[dr] (cstation) -- node[lbl, right, text=rdfsprop, pos=0.45] {rdfs:range} (cs);
  \draw[dr] (bprov.south east) -- node[lbl, sloped, above, text=rdfsprop] {rdfs:range} (prov.north west);

  % ---- instance triples ----
  \draw[e] (cp) -- node[lbl, above] {vio:careerStation} (st2);
  \draw[e] (st2) -- node[lbl, above] {vio:team} (hagl);
  \draw[e] (cp.south) -- ++(0,-0.55) -| node[lbl, pos=0.25] {vio:birthProvince} (na.south);

  % ---- rdf:type across the line ----
  \draw[typ] (cp)   -- node[lbl, right] {rdf:type} (fp);
  \draw[typ] (st2)  -- node[lbl, right] {rdf:type} (club);
  \draw[typ] (hagl) -- node[lbl, right] {rdf:type} (fc);
  \draw[typ] (na)   -- node[lbl, right] {rdf:type} (prov);

  % ---- a triple entailed by RDFS (rdfs9: type propagates up subClassOf) ----
  \draw[->, rdfsaxiom, dashed, thick] (cp.west) to[out=170, in=190, looseness=1.15]
      node[lbl, text=rdfsaxiom, sloped, above, pos=0.5] {rdf:type (entailed)} (dboperson.west);
\end{tikzpicture}
) and slides/
% (% RDFS graph of the vio: ontology in the lecture style: ellipses for resources,
% rdfs:Class / rdf:Property at the top, the schema above a horizontal line and
% instance data below it, rdf:type edges crossing the line.
% An excerpt: the full class hierarchy is in the report's class table.
%
% Shared by report/ (\input{figures/rdfs_graph}) and slides/
% (\input{../report/figures/rdfs_graph}). The including preamble defines the
% colours rdfsdraw, rdfsfill (instances), rdfsprop (properties) and rdfsaxiom
% (entailed triples).
\begin{tikzpicture}[
    % absolute font sizes, so the figure looks the same in the report (10pt) and in beamer (11pt)
    font=\fontsize{7}{8.5}\selectfont, >={Stealth[length=1.7mm]},
    res/.style={ellipse, draw=black!80, fill=white, inner sep=1pt, minimum height=0.72cm,
                font=\fontsize{7}{8.5}\selectfont\ttfamily},
    meta/.style={res, fill=black!8, draw=black!60},
    dbo/.style={res, dashed, draw=black!55, text=black!70},
    prp/.style={res, draw=rdfsprop, text=rdfsprop},
    dboprp/.style={prp, dashed},
    inst/.style={res, fill=rdfsfill, draw=rdfsdraw},
    e/.style={->, black!80},
    typ/.style={->, black!45},
    dr/.style={->, rdfsprop},
    lbl/.style={font=\fontsize{5}{6}\selectfont\itshape, fill=white, inner sep=0.8pt},
  ]
  % ---- RDFS vocabulary ----
  \node[meta] (class)    at (4.7,1.75)  {rdfs:Class};
  \node[meta] (property) at (10.6,1.75) {rdf:Property};

  % ---- classes ----
  \node[dbo] (dboanimal) at (1.8,1.75)  {dbo:Animal};
  \node[dbo] (dboperson) at (1.8,0.35)  {dbo:Person};
  \node[res] (person)    at (1.8,-1.05) {vio:Person};
  \node[res] (athlete)   at (1.8,-2.45) {vio:Athlete};
  \node[res] (fp)        at (1.8,-3.85) {vio:FootballPlayer};
  \node[res] (cs)        at (6.2,-2.45) {vio:CareerStation};
  \node[res] (club)      at (6.2,-3.85) {vio:ClubStation};
  \node[res] (org)       at (10.4,-2.45) {vio:Organisation};
  \node[res] (fc)        at (10.4,-3.85) {vio:FootballClub};
  \node[res] (loc)       at (14.7,-2.45) {vio:Location};
  \node[res] (prov)      at (14.7,-3.85) {vio:Province};

  % ---- properties ----
  \node[prp]    (cstation)  at (6.2,0.35)   {vio:careerStation};
  \node[prp]    (bplace)    at (10.6,0.35)  {vio:birthPlace};
  \node[dboprp] (dbobplace) at (15.0,0.35)  {dbo:birthPlace};
  \node[prp]    (bprov)     at (10.6,-1.05) {vio:birthProvince};

  % ---- instance data ----
  \node[inst] (cp)   at (1.8,-6.0)  {vres:Nguyễn\_Công\_Phượng};
  \node[inst] (st2)  at (6.2,-6.0)  {vres:…\_\_2};
  \node[inst] (hagl) at (10.4,-6.0) {vres:…Hoàng\_Anh\_Gia\_Lai};
  \node[inst] (na)   at (14.7,-6.0) {vres:Nghệ\_An};

  % ---- schema / data line ----
  \draw[black!45, thick] (-0.9,-4.95) -- (16.7,-4.95);
  \node[lbl, fill=none, text=black!55, anchor=south west, inner sep=2pt] at (-0.9,-4.95) {schema (RDFS)};
  \node[lbl, fill=none, text=black!55, anchor=north west, inner sep=2pt] at (-0.9,-4.95) {data (RDF)};

  % ---- rdfs:subClassOf ----
  \draw[e] (dboperson) -- node[lbl, right] {rdfs:subClassOf} (dboanimal);
  \draw[e] (person)  -- node[lbl, right] {rdfs:subClassOf} (dboperson);
  \draw[e] (athlete) -- node[lbl, right] {rdfs:subClassOf} (person);
  \draw[e] (fp)      -- node[lbl, right] {rdfs:subClassOf} (athlete);
  \draw[e] (club)    -- node[lbl, right] {rdfs:subClassOf} (cs);
  \draw[e] (fc)      -- node[lbl, right] {rdfs:subClassOf} (org);
  \draw[e] (prov)    -- node[lbl, right] {rdfs:subClassOf} (loc);

  % ---- rdfs:subPropertyOf ----
  \draw[e] (bprov)  -- node[lbl, right] {rdfs:subPropertyOf} (bplace);
  \draw[e] (bplace) -- node[lbl, above] {rdfs:subPropertyOf} (dbobplace);

  % ---- rdf:type to the RDFS vocabulary ----
  \draw[typ] (person.north east) -- node[lbl, sloped, above] {rdf:type} (class.south west);
  \draw[typ] (cstation) -- node[lbl, sloped, above] {rdf:type} (property);
  \draw[typ] (bplace) -- node[lbl, right] {rdf:type} (property);
  \draw[typ] (dbobplace) -- node[lbl, sloped, above] {rdf:type} (property);

  % ---- rdfs:domain / rdfs:range ----
  \draw[dr] (cstation.south west) -- node[lbl, sloped, above, text=rdfsprop] {rdfs:domain} (fp.north east);
  \draw[dr] (cstation) -- node[lbl, right, text=rdfsprop, pos=0.45] {rdfs:range} (cs);
  \draw[dr] (bprov.south east) -- node[lbl, sloped, above, text=rdfsprop] {rdfs:range} (prov.north west);

  % ---- instance triples ----
  \draw[e] (cp) -- node[lbl, above] {vio:careerStation} (st2);
  \draw[e] (st2) -- node[lbl, above] {vio:team} (hagl);
  \draw[e] (cp.south) -- ++(0,-0.55) -| node[lbl, pos=0.25] {vio:birthProvince} (na.south);

  % ---- rdf:type across the line ----
  \draw[typ] (cp)   -- node[lbl, right] {rdf:type} (fp);
  \draw[typ] (st2)  -- node[lbl, right] {rdf:type} (club);
  \draw[typ] (hagl) -- node[lbl, right] {rdf:type} (fc);
  \draw[typ] (na)   -- node[lbl, right] {rdf:type} (prov);

  % ---- a triple entailed by RDFS (rdfs9: type propagates up subClassOf) ----
  \draw[->, rdfsaxiom, dashed, thick] (cp.west) to[out=170, in=190, looseness=1.15]
      node[lbl, text=rdfsaxiom, sloped, above, pos=0.5] {rdf:type (entailed)} (dboperson.west);
\end{tikzpicture}
). The including preamble defines the
% colours rdfsdraw, rdfsfill (instances), rdfsprop (properties) and rdfsaxiom
% (entailed triples).
\begin{tikzpicture}[
    % absolute font sizes, so the figure looks the same in the report (10pt) and in beamer (11pt)
    font=\fontsize{7}{8.5}\selectfont, >={Stealth[length=1.7mm]},
    res/.style={ellipse, draw=black!80, fill=white, inner sep=1pt, minimum height=0.72cm,
                font=\fontsize{7}{8.5}\selectfont\ttfamily},
    meta/.style={res, fill=black!8, draw=black!60},
    dbo/.style={res, dashed, draw=black!55, text=black!70},
    prp/.style={res, draw=rdfsprop, text=rdfsprop},
    dboprp/.style={prp, dashed},
    inst/.style={res, fill=rdfsfill, draw=rdfsdraw},
    e/.style={->, black!80},
    typ/.style={->, black!45},
    dr/.style={->, rdfsprop},
    lbl/.style={font=\fontsize{5}{6}\selectfont\itshape, fill=white, inner sep=0.8pt},
  ]
  % ---- RDFS vocabulary ----
  \node[meta] (class)    at (4.7,1.75)  {rdfs:Class};
  \node[meta] (property) at (10.6,1.75) {rdf:Property};

  % ---- classes ----
  \node[dbo] (dboanimal) at (1.8,1.75)  {dbo:Animal};
  \node[dbo] (dboperson) at (1.8,0.35)  {dbo:Person};
  \node[res] (person)    at (1.8,-1.05) {vio:Person};
  \node[res] (athlete)   at (1.8,-2.45) {vio:Athlete};
  \node[res] (fp)        at (1.8,-3.85) {vio:FootballPlayer};
  \node[res] (cs)        at (6.2,-2.45) {vio:CareerStation};
  \node[res] (club)      at (6.2,-3.85) {vio:ClubStation};
  \node[res] (org)       at (10.4,-2.45) {vio:Organisation};
  \node[res] (fc)        at (10.4,-3.85) {vio:FootballClub};
  \node[res] (loc)       at (14.7,-2.45) {vio:Location};
  \node[res] (prov)      at (14.7,-3.85) {vio:Province};

  % ---- properties ----
  \node[prp]    (cstation)  at (6.2,0.35)   {vio:careerStation};
  \node[prp]    (bplace)    at (10.6,0.35)  {vio:birthPlace};
  \node[dboprp] (dbobplace) at (15.0,0.35)  {dbo:birthPlace};
  \node[prp]    (bprov)     at (10.6,-1.05) {vio:birthProvince};

  % ---- instance data ----
  \node[inst] (cp)   at (1.8,-6.0)  {vres:Nguyễn\_Công\_Phượng};
  \node[inst] (st2)  at (6.2,-6.0)  {vres:…\_\_2};
  \node[inst] (hagl) at (10.4,-6.0) {vres:…Hoàng\_Anh\_Gia\_Lai};
  \node[inst] (na)   at (14.7,-6.0) {vres:Nghệ\_An};

  % ---- schema / data line ----
  \draw[black!45, thick] (-0.9,-4.95) -- (16.7,-4.95);
  \node[lbl, fill=none, text=black!55, anchor=south west, inner sep=2pt] at (-0.9,-4.95) {schema (RDFS)};
  \node[lbl, fill=none, text=black!55, anchor=north west, inner sep=2pt] at (-0.9,-4.95) {data (RDF)};

  % ---- rdfs:subClassOf ----
  \draw[e] (dboperson) -- node[lbl, right] {rdfs:subClassOf} (dboanimal);
  \draw[e] (person)  -- node[lbl, right] {rdfs:subClassOf} (dboperson);
  \draw[e] (athlete) -- node[lbl, right] {rdfs:subClassOf} (person);
  \draw[e] (fp)      -- node[lbl, right] {rdfs:subClassOf} (athlete);
  \draw[e] (club)    -- node[lbl, right] {rdfs:subClassOf} (cs);
  \draw[e] (fc)      -- node[lbl, right] {rdfs:subClassOf} (org);
  \draw[e] (prov)    -- node[lbl, right] {rdfs:subClassOf} (loc);

  % ---- rdfs:subPropertyOf ----
  \draw[e] (bprov)  -- node[lbl, right] {rdfs:subPropertyOf} (bplace);
  \draw[e] (bplace) -- node[lbl, above] {rdfs:subPropertyOf} (dbobplace);

  % ---- rdf:type to the RDFS vocabulary ----
  \draw[typ] (person.north east) -- node[lbl, sloped, above] {rdf:type} (class.south west);
  \draw[typ] (cstation) -- node[lbl, sloped, above] {rdf:type} (property);
  \draw[typ] (bplace) -- node[lbl, right] {rdf:type} (property);
  \draw[typ] (dbobplace) -- node[lbl, sloped, above] {rdf:type} (property);

  % ---- rdfs:domain / rdfs:range ----
  \draw[dr] (cstation.south west) -- node[lbl, sloped, above, text=rdfsprop] {rdfs:domain} (fp.north east);
  \draw[dr] (cstation) -- node[lbl, right, text=rdfsprop, pos=0.45] {rdfs:range} (cs);
  \draw[dr] (bprov.south east) -- node[lbl, sloped, above, text=rdfsprop] {rdfs:range} (prov.north west);

  % ---- instance triples ----
  \draw[e] (cp) -- node[lbl, above] {vio:careerStation} (st2);
  \draw[e] (st2) -- node[lbl, above] {vio:team} (hagl);
  \draw[e] (cp.south) -- ++(0,-0.55) -| node[lbl, pos=0.25] {vio:birthProvince} (na.south);

  % ---- rdf:type across the line ----
  \draw[typ] (cp)   -- node[lbl, right] {rdf:type} (fp);
  \draw[typ] (st2)  -- node[lbl, right] {rdf:type} (club);
  \draw[typ] (hagl) -- node[lbl, right] {rdf:type} (fc);
  \draw[typ] (na)   -- node[lbl, right] {rdf:type} (prov);

  % ---- a triple entailed by RDFS (rdfs9: type propagates up subClassOf) ----
  \draw[->, rdfsaxiom, dashed, thick] (cp.west) to[out=170, in=190, looseness=1.15]
      node[lbl, text=rdfsaxiom, sloped, above, pos=0.5] {rdf:type (entailed)} (dboperson.west);
\end{tikzpicture}
). The including preamble defines the
% colours rdfsdraw, rdfsfill (instances), rdfsprop (properties) and rdfsaxiom
% (entailed triples).
\begin{tikzpicture}[
    % absolute font sizes, so the figure looks the same in the report (10pt) and in beamer (11pt)
    font=\fontsize{7}{8.5}\selectfont, >={Stealth[length=1.7mm]},
    res/.style={ellipse, draw=black!80, fill=white, inner sep=1pt, minimum height=0.72cm,
                font=\fontsize{7}{8.5}\selectfont\ttfamily},
    meta/.style={res, fill=black!8, draw=black!60},
    dbo/.style={res, dashed, draw=black!55, text=black!70},
    prp/.style={res, draw=rdfsprop, text=rdfsprop},
    dboprp/.style={prp, dashed},
    inst/.style={res, fill=rdfsfill, draw=rdfsdraw},
    e/.style={->, black!80},
    typ/.style={->, black!45},
    dr/.style={->, rdfsprop},
    lbl/.style={font=\fontsize{5}{6}\selectfont\itshape, fill=white, inner sep=0.8pt},
  ]
  % ---- RDFS vocabulary ----
  \node[meta] (class)    at (4.7,1.75)  {rdfs:Class};
  \node[meta] (property) at (10.6,1.75) {rdf:Property};

  % ---- classes ----
  \node[dbo] (dboanimal) at (1.8,1.75)  {dbo:Animal};
  \node[dbo] (dboperson) at (1.8,0.35)  {dbo:Person};
  \node[res] (person)    at (1.8,-1.05) {vio:Person};
  \node[res] (athlete)   at (1.8,-2.45) {vio:Athlete};
  \node[res] (fp)        at (1.8,-3.85) {vio:FootballPlayer};
  \node[res] (cs)        at (6.2,-2.45) {vio:CareerStation};
  \node[res] (club)      at (6.2,-3.85) {vio:ClubStation};
  \node[res] (org)       at (10.4,-2.45) {vio:Organisation};
  \node[res] (fc)        at (10.4,-3.85) {vio:FootballClub};
  \node[res] (loc)       at (14.7,-2.45) {vio:Location};
  \node[res] (prov)      at (14.7,-3.85) {vio:Province};

  % ---- properties ----
  \node[prp]    (cstation)  at (6.2,0.35)   {vio:careerStation};
  \node[prp]    (bplace)    at (10.6,0.35)  {vio:birthPlace};
  \node[dboprp] (dbobplace) at (15.0,0.35)  {dbo:birthPlace};
  \node[prp]    (bprov)     at (10.6,-1.05) {vio:birthProvince};

  % ---- instance data ----
  \node[inst] (cp)   at (1.8,-6.0)  {vres:Nguyễn\_Công\_Phượng};
  \node[inst] (st2)  at (6.2,-6.0)  {vres:…\_\_2};
  \node[inst] (hagl) at (10.4,-6.0) {vres:…Hoàng\_Anh\_Gia\_Lai};
  \node[inst] (na)   at (14.7,-6.0) {vres:Nghệ\_An};

  % ---- schema / data line ----
  \draw[black!45, thick] (-0.9,-4.95) -- (16.7,-4.95);
  \node[lbl, fill=none, text=black!55, anchor=south west, inner sep=2pt] at (-0.9,-4.95) {schema (RDFS)};
  \node[lbl, fill=none, text=black!55, anchor=north west, inner sep=2pt] at (-0.9,-4.95) {data (RDF)};

  % ---- rdfs:subClassOf ----
  \draw[e] (dboperson) -- node[lbl, right] {rdfs:subClassOf} (dboanimal);
  \draw[e] (person)  -- node[lbl, right] {rdfs:subClassOf} (dboperson);
  \draw[e] (athlete) -- node[lbl, right] {rdfs:subClassOf} (person);
  \draw[e] (fp)      -- node[lbl, right] {rdfs:subClassOf} (athlete);
  \draw[e] (club)    -- node[lbl, right] {rdfs:subClassOf} (cs);
  \draw[e] (fc)      -- node[lbl, right] {rdfs:subClassOf} (org);
  \draw[e] (prov)    -- node[lbl, right] {rdfs:subClassOf} (loc);

  % ---- rdfs:subPropertyOf ----
  \draw[e] (bprov)  -- node[lbl, right] {rdfs:subPropertyOf} (bplace);
  \draw[e] (bplace) -- node[lbl, above] {rdfs:subPropertyOf} (dbobplace);

  % ---- rdf:type to the RDFS vocabulary ----
  \draw[typ] (person.north east) -- node[lbl, sloped, above] {rdf:type} (class.south west);
  \draw[typ] (cstation) -- node[lbl, sloped, above] {rdf:type} (property);
  \draw[typ] (bplace) -- node[lbl, right] {rdf:type} (property);
  \draw[typ] (dbobplace) -- node[lbl, sloped, above] {rdf:type} (property);

  % ---- rdfs:domain / rdfs:range ----
  \draw[dr] (cstation.south west) -- node[lbl, sloped, above, text=rdfsprop] {rdfs:domain} (fp.north east);
  \draw[dr] (cstation) -- node[lbl, right, text=rdfsprop, pos=0.45] {rdfs:range} (cs);
  \draw[dr] (bprov.south east) -- node[lbl, sloped, above, text=rdfsprop] {rdfs:range} (prov.north west);

  % ---- instance triples ----
  \draw[e] (cp) -- node[lbl, above] {vio:careerStation} (st2);
  \draw[e] (st2) -- node[lbl, above] {vio:team} (hagl);
  \draw[e] (cp.south) -- ++(0,-0.55) -| node[lbl, pos=0.25] {vio:birthProvince} (na.south);

  % ---- rdf:type across the line ----
  \draw[typ] (cp)   -- node[lbl, right] {rdf:type} (fp);
  \draw[typ] (st2)  -- node[lbl, right] {rdf:type} (club);
  \draw[typ] (hagl) -- node[lbl, right] {rdf:type} (fc);
  \draw[typ] (na)   -- node[lbl, right] {rdf:type} (prov);

  % ---- a triple entailed by RDFS (rdfs9: type propagates up subClassOf) ----
  \draw[->, rdfsaxiom, dashed, thick] (cp.west) to[out=170, in=190, looseness=1.15]
      node[lbl, text=rdfsaxiom, sloped, above, pos=0.5] {rdf:type (entailed)} (dboperson.west);
\end{tikzpicture}
}
\caption{RDFS graph of an excerpt of the ontology. Above the line is the
schema: classes and properties are resources typed by \term{rdfs:Class} and
\term{rdf:Property}, linked by \term{rdfs:subClassOf} and
\term{rdfs:subPropertyOf}; blue edges are the \term{rdfs:domain} and
\term{rdfs:range} of two properties. Below the line are instance data
(\term{vres:…\_\_2} is the career station \term{vres:Nguyễn\_Công\_Phượng\_\_2}).
The dashed edge is a triple that RDFS entails: the player is also a
\term{dbo:Person}. \term{dbo:Person}~$\sqsubseteq$~\term{dbo:Animal} is DBpedia's own axiom. Only some \term{rdf:type} edges to the RDFS vocabulary are
drawn; the full class hierarchy is in Table~\ref{tab:classes}.}
\label{fig:rdfs}
\end{figure}

\begin{table}[t]
\centering
\caption{Class hierarchy of \term{vio:} (indentation shows
\term{rdfs:subClassOf}), the DBpedia class each class is aligned to, and the
number of instances before and after OWL~2~RL reasoning.}
\label{tab:classes}
\small
\begin{adjustbox}{max width=\textwidth}
\begin{tabular}{@{}llrr@{}}
\toprule
Class & DBpedia superclass & Asserted & After reasoning \\
\midrule
\term{Person}                             & \term{dbo:Person}                 & 0     & 659 \\
\quad\term{Athlete}                       & \term{dbo:Athlete}                & 0     & 611 \\
\qquad\term{FootballPlayer}               & \term{dbo:SoccerPlayer}           & 611   & 611 \\
\qquad\quad\term{NationalTeamPlayer}      & ---                               & 0     & 486 \\
\term{Organisation}                       & \term{dbo:Organisation}           & 0     & 410 \\
\quad\term{FootballClub}                  & \term{dbo:SoccerClub}             & 61    & 215 \\
\quad\term{NationalFootballTeam}          & \term{dbo:NationalSoccerClub}     & 10    & 27 \\
\quad\term{EducationalInstitution}        & \term{dbo:EducationalInstitution} & 0     & 168 \\
\qquad\term{University}                   & \term{dbo:University}             & 168   & 168 \\
\term{Location}                           & \term{dbo:Place}                  & 0     & 370 \\
\quad\term{Province}                      & \term{dbo:Province}               & 110   & 110 \\
\qquad\term{FormerProvince}               & ---                               & 76    & 76 \\
\quad\term{Stadium}                       & \term{dbo:Stadium}                & 29    & 44 \\
\quad\term{Country}                       & \term{dbo:Country}                & 1     & 1 \\
\term{CareerStation}                      & \term{dbo:CareerStation}          & 0     & 3{,}382 \\
\quad\term{YouthStation}                  & ---                               & 527     & 527 \\
\quad\term{ClubStation}                   & ---                               & 1{,}966 & 1{,}966 \\
\quad\term{NationalTeamStation}           & ---                               & 889     & 889 \\
\bottomrule
\end{tabular}
\end{adjustbox}
\end{table}

\paragraph{Properties.} The ontology has 54 properties, 26 object and 28
datatype properties (Figure~\ref{fig:rdfs}). All have labels in Vietnamese and
English. Every datatype property has an \term{rdfs:range}, and 39 properties
have an \term{rdfs:domain} (24 datatype, 15 object). When DBpedia has an
equivalent of the same kind, the property is an \term{rdfs:subPropertyOf} of it
(38 properties). Typical examples are
\term{vio:birthDate} $\sqsubseteq$ \term{dbo:birthDate},
\term{vio:capacity} $\sqsubseteq$ \term{dbo:seatingCapacity}, and
\term{vio:birthProvince} $\sqsubseteq$ \term{vio:birthPlace} $\sqsubseteq$
\term{dbo:birthPlace}. 22 datatype properties are functional, for example
birth date, height and capacity. Some domains are left open on purpose. If
\term{vio:latitude} or \term{vio:province} had the domain \term{Location},
every club and university that has coordinates or a province would be
inferred to be a location, which contradicts the disjointness axioms.

\paragraph{Axioms.} A career is modelled with the ``role at a time'' pattern
of \term{dbo:CareerStation}. A player has one \term{vio:careerStation} per
team and period, and the station carries the team, the start and end years,
appearances, goals and an on-loan flag. Four kinds of axiom act on this
structure during reasoning:
\begin{itemize}
  \item a \emph{property chain},
        $\term{vio:careerStation} \circ \term{vio:team} \sqsubseteq \term{vio:playedFor}$,
        which gives a direct player--team edge;
  \item \emph{inverses}, so that each relation can be queried in both
        directions: \term{vio:hasPlayer} is the inverse of
        \term{vio:playedFor}, and the same holds for the pairs
        \term{vio:hasPart}/\term{vio:isPartOf} and
        \term{vio:predecessor}/\term{vio:successor};
  \item \emph{restrictions}, which derive classes from the career data:
        {\small
        \begin{gather*}
          \term{FootballPlayer} \sqcap \exists\,\term{careerStation}.\term{NationalTeamStation} \sqsubseteq \term{NationalTeamPlayer},\\
          \term{ClubStation} \sqsubseteq \forall\,\term{team}.\term{FootballClub}.
        \end{gather*}}
        The intersection keeps the first definition independent of the domain
        of \term{careerStation}~\cite{paulheim2024modeling}. Similar $\forall$
        restrictions hold for youth and national-team stations. They give a
        class to teams that are not seeds, such as the Japanese club Mito
        HollyHock;
  \item four \emph{disjointness} axioms (\term{owl:AllDisjointClasses}): between
        the four roots, between club, national team and educational
        institution, between province, stadium and country, and between the
        three kinds of station. They let the reasoner detect modelling errors
        (Section~\ref{sec:reasoning}).
\end{itemize}

% =========================================================================
\subsection{Data Collection}
\label{sec:collection}

\paragraph{Entity selection.} \term{seeds} runs one Wikidata Query Service
query per class. A query matches the class through \term{P31}/\term{P279*} and
the country through \term{P17} or \term{P27} (Vietnam, \term{Q881}). It keeps
only items with a Vietnamese Wikipedia article, and makes the English article
optional. The query for players is shown below. Wikimedia list pages,
disambiguation pages and titles starting with ``Danh sách'' (list of) are
removed. When an item matches several classes, the first class in the
priority order wins.
\begin{Verbatim}
SELECT DISTINCT ?item ?viTitle ?enTitle WHERE {
  ?item wdt:P31 wd:Q5 ; wdt:P106 wd:Q937857 ; wdt:P27 wd:Q881 .  # footballer, Vietnam
  FILTER NOT EXISTS { ?item wdt:P31 wd:Q13406463 }               # Wikimedia list page
  FILTER NOT EXISTS { ?item wdt:P31 wd:Q4167410 }                # disambiguation page
  ?vs schema:about ?item ; schema:isPartOf <https://vi.wikipedia.org/> ;
      schema:name ?viTitle .
  OPTIONAL { ?es schema:about ?item ; schema:isPartOf <https://en.wikipedia.org/> ;
                 schema:name ?enTitle . } }
\end{Verbatim}
Facts are fetched in blocks of at most 300 items: dates with their precision,
so that a year-only birth date becomes \term{xsd:gYear} instead of a fabricated
1~January; quantities (height, capacity, population, area, students) with their
point in time (\term{P585}), so that the latest value is kept; relations with
their start, end, appearance and goal qualifiers; and the labels and sitelinks
of every referenced item. Birthplaces and locations are mapped to a province
through \term{P131*}. The result is 990 seeds: 611 players, 61 clubs, 10
national teams, 29 stadiums, 168 universities, 110 provinces and 1 country.

\paragraph{Article collection.} \term{enrich} reads the Vietnamese articles
through the MediaWiki Action API, in batches of up to 50 titles with
\term{maxlag} set. It stores the page ID, revision, length, coordinates, lead
image, plain-text introduction (the abstract), non-hidden categories and
incoming redirects, and, instead of the whole wikitext, the parameters of the
outermost infobox (found with \term{mwparserfromhell}) and the URLs of the
``Liên kết ngoài'' (external links) section. Every title linked from an infobox
is resolved, following redirects and flagging disambiguation pages (2{,}588
titles, 71 disambiguation pages), and redirects to the infobox templates are
kept as template aliases. Requests are spaced 2\,s apart for the Wikidata Query
Service and 1\,s for Wikipedia, carry a contact User-Agent, and are retried
on HTTP 429 and \term{maxlag}; \term{-{}-offline} and
\term{-{}-refresh} use only the cache or bypass it.

% =========================================================================
\subsection{Transformation to RDF}
\label{sec:transform}

\paragraph{IRIs.} A resource IRI is \term{vres:} followed by the Vietnamese
title as MediaWiki normalises it: NFC Unicode, first character in upper case,
spaces as underscores. Like DBpedia chapters, the IRI keeps Unicode characters,
for example \term{vres:Nguyễn\_Công\_Phượng}, and percent-encodes only the
characters that are not allowed in an IRI. Career stations are numbered from
youth to club to national team (\term{vres:Nguyễn\_Công\_Phượng\_\_3}).
Redirects are separate resources that point to the main resource with
\term{dbo:wikiPageRedirects}, and their titles are also added to the main
resource as \term{skos:altLabel}, so that it can be found by its other names.

\paragraph{DBpedia datasets.} For every entity, \term{build} produces the
same basic datasets as DBpedia (Table~\ref{tab:core}). Raw infobox
values are kept under \term{vip:}, as in DBpedia's \term{dbp:} namespace.

\begin{table}[t]
\centering
\caption{DBpedia datasets reproduced for every entity, with their size in the
current build.}
\label{tab:core}
\small
\begin{tabularx}{\textwidth}{@{}l>{\raggedright\arraybackslash}Xr@{}}
\toprule
Dataset & Properties & Size \\
\midrule
Labels & \term{rdfs:label} @vi, @en & 990 vi, 805 en \\
Abstracts & \term{dbo:abstract} (lead section), \term{rdfs:comment} (at most 400 characters) & 990 \\
Page metadata & \term{dbo:wikiPageID}, \term{dbo:wikiPageRevisionID}, \term{prov:wasDerivedFrom}, \term{foaf:isPrimaryTopicOf} & 990 \\
Images & \term{dbo:thumbnail}, \term{foaf:depiction} & 543 \\
Categories & \term{dct:subject} $\rightarrow$ \term{skos:Concept} & 9{,}126 links, 1{,}154 categories \\
Redirects & \term{dbo:wikiPageRedirects}, \term{skos:altLabel} & 1{,}274 \\
External links & \term{dbo:wikiPageExternalLink}, \term{foaf:homepage} & 713 links, 270 homepages \\
Coordinates & \term{vio:latitude}/\term{vio:longitude} ($\sqsubseteq$ \term{wgs84:}), \term{georss:point} & 189 entities \\
Interlinks & \term{owl:sameAs} \term{dbr:}, \term{wd:} & 777, 990 \\
Raw infobox & \term{vip:}\emph{key} & 24{,}605 triples, 687 properties \\
\bottomrule
\end{tabularx}
\end{table}

\paragraph{Mapping-based infobox extraction.} Infobox values are mapped to
\term{vio:} properties in four stages. The wikitext of each value is first
\emph{rendered}: common nested templates (dates, heights, URLs and the list
templates \term{ubl}, \term{plainlist} and \term{hlist}) are interpreted, and
references, comments and flag icons removed. It is then \emph{parsed} into a
typed literal: the parsers understand Vietnamese numbers (\term{16.361,2}),
open year ranges (\term{2015–}) and loan markers (\term{→}, \term{(mượn)}), and
accept heights only between 1.4 and 2.2\,m. Keys are \emph{mapped} to
\term{vio:} properties by the \term{MAPPINGS} table of five template families
(players, clubs and national teams, stadiums, universities, provinces); as in
the mappings wiki of DBpedia, a property can take keys from both the Vietnamese
and the English template. Finally the \term{PRECEDENCE} table \emph{chooses a
source} when Wikidata and the infobox both have a value: Wikidata wins for
dates, height, population, area, capacity, opening year, homepage and
abbreviation, the infobox for shirt number, numbers of students and academic
staff, motto and administrative code, and object values of both are combined.
Link targets are filtered by class: the team of a career station cannot be a
province, a stadium or a university, and the \term{vio:ground} of a club must
be a stadium.

Career stations come from the numbered parameters of the football biography
infobox (\term{years}$n$, \term{clubs}$n$, \term{caps}$n$ and \term{goals}$n$)
and their \term{youth} and \term{national} variants. For 594 of the 611 players
the career is read from the infobox; the other 17 have no club rows (or no
infobox) and fall back to the Wikidata statements \term{P54}. In total 601
players receive 3{,}382 stations. The
build also writes \term{mapping\_stats.json}, which counts mapped
and unmapped keys for each template family, to guide which mappings to add next.

\paragraph{From three to four stars.} Listing~\ref{lst:fourstar} follows one
infobox through the transformation. Before it, the data are machine-readable and
open but name nothing with a URI: at most three stars. After it, they are four-star RDF, through four changes.
\emph{Every thing gets an HTTP IRI}: the player, the district, the province,
both clubs and even the loan spell, \term{vres:Nguyễn\_Công\_Phượng\_\_3}.
\emph{Values become typed literals} (\term{xsd:date}, \term{xsd:gYear},
\term{xsd:nonNegativeInteger}, \term{xsd:double}, \term{xsd:boolean}) and texts
carry a language tag. \emph{Markup is interpreted}: templates such as
\term{ngày sinh và tuổi} (date of birth and age) become values, links are
resolved through redirects, the loan marker \term{→ … (mượn)} becomes
\term{vio:onLoan true}, and \term{PRECEDENCE} takes the height of 1.68\,m from
Wikidata over the infobox's 1.69\,m. \emph{Terms come from shared
vocabularies}: \term{vio:} terms specialise \term{dbo:} terms, and the core
datasets use \term{rdfs:}, \term{skos:}, \term{foaf:}, \term{dct:} and
\term{prov:}.
\begin{listing}[t]
\begin{Verbatim}
{{Thông tin tiểu sử bóng đá
| birth_date  = {{ngày sinh và tuổi|1995|1|21}}<ref name=a/>
| birth_place = [[Đô Lương (huyện)|Đô Lương]], [[Nghệ An]], [[Việt Nam]]
| height      = {{convert|1,69|m}}
| currentclub = [[Câu lạc bộ bóng đá Thành phố Đồng Nai|Thành phố Đồng Nai]]
| clubnumber  = 10
| years2 = 2016   | clubs2 = → [[Mito HollyHock]] (mượn)   | caps2 = 5   | goals2 = 0
\end{Verbatim}
\vspace{-0.4em}
\begin{center}\footnotesize $\Downarrow$ \quad \texttt{python -m vidbpedia build}\end{center}
\vspace{-0.6em}
\begin{Verbatim}
vres:Nguyễn_Công_Phượng a vio:FootballPlayer ;
    vio:birthDate "1995-01-21"^^xsd:date ;
    vio:birthPlace vres:Đô_Lương_\(huyện\) ;
    vio:birthProvince vres:Nghệ_An ;                   # Wikidata P19, then P131*
    vio:height 1.68e+00 ;                              # xsd:double, from Wikidata
    vio:currentClub vres:Câu_lạc_bộ_bóng_đá_Thành_phố_Đồng_Nai ;
    vio:shirtNumber "10"^^xsd:nonNegativeInteger ;
    vio:careerStation vres:Nguyễn_Công_Phượng__3 .
vres:Nguyễn_Công_Phượng__3 a vio:ClubStation ;
    vio:team vres:Mito_HollyHock ;  vio:onLoan true ;
    vio:startYear "2016"^^xsd:gYear ;  vio:endYear "2016"^^xsd:gYear ;
    vio:appearances "5"^^xsd:nonNegativeInteger ;  vio:goals "0"^^xsd:nonNegativeInteger .
\end{Verbatim}
\caption{From three to four stars: part of the infobox of Nguyễn Công Phượng
(top, wikitext) and the triples the build produces for it (bottom, Turtle,
excerpt).}
\label{lst:fourstar}
\end{listing}

% =========================================================================
\subsection{Interlinking}
\label{sec:interlinking}

The English link of an entity comes from its Wikidata item: if the item has an
English sitelink, the resource receives \term{owl:sameAs} to the \term{dbr:}
resource of that title, written in the same Unicode IRI form that DBpedia
uses. Using the sitelink of the same item avoids matching by label, which
would be ambiguous for common Vietnamese names. Every entity also receives
\term{owl:sameAs} to its Wikidata item. In addition to these instance links, the ontology is linked
to DBpedia at the schema level through its subclass and sub-property axioms.

The dataset describes itself with VoID~\cite{alexander2011void}: a
\term{void:Dataset} with the licence (CC~BY-SA~4.0, as Wikipedia), sources, URI
space, vocabularies and dump location; three subsets by provenance (ontology,
asserted, inferred) with their triple counts; a class partition for each of the
seven entity classes; and two \term{void:Linkset}s, to DBpedia and to Wikidata,
with the number of links.

% =========================================================================
\subsection{Validation, Reasoning and Export}
\label{sec:postprocess}

\paragraph{Validation.} \term{postprocess} runs eight checks before reasoning
(Table~\ref{tab:checks}). Each reports issues with severity, subject and
detail; the tests run the same functions, and with
\term{-{}-strict} any error stops the build.

\begin{table}[t]
\centering
\caption{Quality checks (\term{kg/validation.py}). The last build reports no
error and no warning.}
\label{tab:checks}
\small
\begin{tabularx}{\textwidth}{@{}l>{\raggedright\arraybackslash}X@{}}
\toprule
Check & Condition \\
\midrule
\term{tbox} & sub-properties link properties of the same kind; every \term{vio:} class has a \term{dbo:} ancestor; terms used in axioms are declared; declared terms have vi and en labels and \term{rdfs:isDefinedBy}; \term{vio:} names are CamelCase \\
\term{property\_kinds} & object properties point to IRIs, datatype properties to literals \\
\term{datatype\_ranges} & literals match the declared \term{rdfs:range} (including \term{rdf:langString}) and are well-formed \\
\term{required} & every entity has exactly one main class, a Vietnamese label, a source page and a Wikidata link \\
\term{functional} & functional properties have at most one value \\
\term{object\_ranges} & if the target of an object property has a class, that class fits the range \\
\term{sameas} & no list or disambiguation targets; no two resources share a target \\
\term{plausibility} & coordinates inside Vietnam; birth year, capacity, population and height in a plausible range (warnings) \\
\bottomrule
\end{tabularx}
\end{table}

\paragraph{Reasoning.}
\label{sec:reasoning}
Reasoning uses the OWL~2~RL rules of the \term{owlrl}
library~\cite{motik2012owl2profiles, rdflib} and runs in 35\,s. Three choices
keep it fast and correct. Functional-property, cardinality and
\term{owl:sameAs} axioms are removed from the TBox, because OWL~2~RL would
otherwise derive \term{x owl:sameAs x} for every node and merge entities that
share a functional value; functional properties are checked by validation
instead. Only the \emph{skeleton} of the data reaches the reasoner, the
\term{rdf:type} triples and the \term{vio:} triples of \term{vres:} resources,
since abstracts, images and raw infobox triples cannot trigger any rule. Only
new triples whose class or property is in \term{vio:}, \term{dbo:} or
\term{wgs84:} are kept, never \term{owl:sameAs}, and they are written to their
own file, \term{inferred.nt}, so the interface can label each triple as
asserted or inferred.
Reasoning also helped us find two modelling errors. First, it reported a
violation of the disjointness between clubs and national teams. The cause was
that \term{vio:ground}, \term{vio:manager} and \term{vio:chairman} had the
domain \term{FootballClub}, so every national team with a home ground was
inferred to be a club; the domain is now \term{Organisation}. Second, the
question ``which clubs have their home ground in Hà Nội'' returned clubs from
Đà Nẵng and Khánh Hòa. The cause was an axiom declaring \term{vio:tenant} the
inverse of \term{vio:ground}, while the ``tenants'' field of the Hàng Đẫy
stadium lists every team that ever played there. We removed the axiom, as
DBpedia also keeps \term{dbo:tenant} and \term{dbo:ground} separate, and kept
only tenants that are entities of the dataset. Figure~\ref{fig:example} shows
asserted and inferred triples around one player.

\paragraph{Export.} The full dataset (ontology, asserted and inferred triples,
VoID) is written as Turtle and N-Triples, and as RDF/XML with
\term{-{}-rdfxml}; the three parts are also written separately. The N-Triples
file is read back and its triple count compared with the graph in memory, and
statistics and the validation report are written as JSON.

\begin{figure}[t]
\centering
\begin{adjustbox}{max width=\textwidth}
% Một thực thể dưới dạng đồ thị RDF theo ký hiệu bài giảng (02_RDF tr.10-14, 04_LOD tr.37):
% elip = tài nguyên (cá thể, lớp), chữ nhật = literal, mỗi cạnh là một triple ghi qname.
% Sinh từ entity_graph_spec.py; slide 10 của bản editable dùng cùng toạ độ.
\begin{tikzpicture}[
    font=\fontsize{7}{8.5}\selectfont, >={Stealth[length=1.7mm]},
    base/.style={inner sep=1.2pt, font=\fontsize{7}{8.5}\selectfont\ttfamily},
    inst/.style={base, ellipse, draw=rdfsdraw, fill=rdfsfill, minimum height=0.7cm},
    cls/.style={base, ellipse, draw=black!80, fill=white, minimum height=0.7cm},
    dbocls/.style={cls, dashed, draw=black!55, text=black!70},
    lod/.style={base, ellipse, draw=accblue, dashed, text=accblue!80!black, minimum height=0.7cm},
    lit/.style={base, rectangle, draw=black!70, fill=black!4, minimum height=0.5cm, inner xsep=3pt},
    e/.style={->, black!80},
    typ/.style={->, black!45},
    inf/.style={->, rdfsaxiom, dashed, thick},
    same/.style={->, accblue, thick},
    lbl/.style={font=\fontsize{5.5}{6.5}\selectfont\itshape, fill=white, inner sep=0.8pt},
  ]
  \node[inst] (cp) at (2.0,0.0) {vres:Nguyễn\_Công\_Phượng};
  \node[inst] (s3) at (7.4,0.0) {vres:…\_\_3};
  \node[inst] (mito) at (11.6,0.0) {vres:Mito\_HollyHock};
  \node[cls] (fp) at (0.0,2.3) {vio:FootballPlayer};
  \node[dbocls] (sp) at (3.9,2.3) {dbo:SoccerPlayer};
  \node[cls] (cstat) at (7.4,2.3) {vio:ClubStation};
  \node[cls] (fc) at (11.6,2.3) {vio:FootballClub};
  \node[lit] (lbl) at (-3.0,0.55) {"Nguyễn Công Phượng"@vi};
  \node[lit] (bd) at (-3.0,-0.55) {"1995-01-21"\textasciicircum{}\textasciicircum{}xsd:date};
  \node[inst] (na) at (-2.6,-2.3) {vres:Nghệ\_An};
  \node[lod] (dbr) at (1.6,-2.3) {dbr:Nguyễn\_Công\_Phượng};
  \node[lod] (wd) at (5.6,-2.3) {wd:Q18045362};
  \node[lit] (sy) at (6.4,-1.45) {"2016"\textasciicircum{}\textasciicircum{}xsd:gYear};
  \node[lit] (ol) at (9.9,-1.3) {true};
  \node[lit] (ap) at (9.4,-2.45) {"5"\textasciicircum{}\textasciicircum{}xsd:nonNegativeInteger};
  \draw[typ] (cp) -- node[lbl] {rdf:type} (fp);
  \draw[inf] (cp) -- node[lbl, text=rdfsaxiom] {rdf:type} (sp);
  \draw[typ] (s3) -- node[lbl] {rdf:type} (cstat);
  \draw[inf] (mito) -- node[lbl, text=rdfsaxiom] {rdf:type} (fc);
  \draw[e] (cp) -- node[lbl] {vio:careerStation} (s3);
  \draw[e] (s3) -- node[lbl] {vio:team} (mito);
  \draw[inf] (cp.north east) -- ++(0,1.15) -| node[lbl, text=rdfsaxiom, pos=0.22] {vio:playedFor} (mito.north west);
  \draw[e] (cp) -- node[lbl] {rdfs:label} (lbl);
  \draw[e] (cp) -- node[lbl] {vio:birthDate} (bd);
  \draw[e] (cp) -- node[lbl] {vio:birthProvince} (na);
  \draw[same] (cp) -- node[lbl, text=accblue] {owl:sameAs} (dbr);
  \draw[same] (cp) -- node[lbl, text=accblue] {owl:sameAs} (wd);
  \draw[e] (s3) -- node[lbl] {vio:startYear} (sy);
  \draw[e] (s3) -- node[lbl] {vio:onLoan} (ol);
  \draw[e] (s3) -- node[lbl] {vio:appearances} (ap);
\end{tikzpicture}

\end{adjustbox}
\caption{Nguyễn Công Phượng and one of his eleven career stations as an RDF graph,
drawn as in the lectures: ellipses are resources (entities and classes), rectangles
are literals, and each edge is one triple labelled with its property. Gray and
black edges are asserted. Orange dashed edges are inferred by OWL~2~RL:
\term{vio:playedFor} through the property chain, \term{dbo:SoccerPlayer} through the
subclass axioms, and the class of Mito HollyHock through the $\forall$ restriction of
\term{vio:ClubStation}. Blue edges are \term{owl:sameAs} links to other datasets.
\term{vres:…\_\_3} is \term{vres:Nguyễn\_Công\_Phượng\_\_3}, the loan spell of
Listing~\ref{lst:fourstar}.}
\label{fig:example}
\end{figure}
